% Generated draft: requires mathematical and source review.
\documentclass[11pt]{article}
\usepackage[margin=25mm]{geometry}
\usepackage{fontspec}
\setmainfont{Latin Modern Roman}
\usepackage{amsmath,amssymb,mathtools,unicode-math}
\setmathfont{Latin Modern Math}
\usepackage{enumitem,needspace,xurl,hyperref}
\setlength{\parindent}{0pt}
\setlength{\parskip}{5pt}
\setlength{\emergencystretch}{3em}
\newcommand{\N}{\mathbb{N}}
\newcommand{\Z}{\mathbb{Z}}
\newcommand{\Q}{\mathbb{Q}}
\newcommand{\R}{\mathbb{R}}
\newcommand{\C}{\mathbb{C}}
\newcommand{\F}{\mathbb{F}}
\newcommand{\A}{\mathbb{A}}
\newcommand{\PP}{\mathbb{P}}
\newcommand{\E}{\mathbb{E}}
\newcommand{\eps}{\varepsilon}
\newcommand{\dd}{\,\mathrm{d}}
\begin{document}
\section*{UnsolvedMath ID: 2307031}
\textbf{Source problem code:} catalogue entry
\subsection{Definitions and mathematical statement}
Let \((a_n)_{n\ge1}\) be real numbers satisfying \(a_1>0\) and \(0\le a_n\le n\), and define
\[
b_n=\sum_{\nu=1}^na_\nu,\qquad
c_n=\sum_{\nu=1}^nb_\nu=\sum_{\nu=1}^n(n-\nu+1)a_\nu.
\]
Thus \(b_n,c_n>0\). The source records that
\[
\sum_{n\ge1}(a_n/c_n)^\alpha<\infty\qquad(\alpha>\tfrac12).
\]
Hayman asks which other functions \(f\) guarantee convergence of \(\sum_{n\ge1}f(c_n/a_n)\) for every sequence satisfying the displayed assumptions, and suggests a relation with \(\sum_{n\ge1}f(n^2)\), under unspecified smoothness assumptions on \(f\).

Zero terms are allowed: when \(a_n=0\), \(c_n/a_n\) is not a finite argument. A possible convention is \(f(\infty)=0\), as appropriate to the recorded power functions \(f(t)=t^{-\alpha}\), and then to write
\[
\sum_{n:a_n>0}f(c_n/a_n).
\]
That convention is not provided by the source. The class, sign, domain, decay, and smoothness of \(f\) must also be fixed. The phrase ``converges with'' may mean that convergence of \(\sum f(n^2)\) implies the universal convergence property, or that it characterizes that universal property. It cannot automatically be read as equivalence of convergence for each individual sequence: the permitted finitely supported sequences would give a finite left sum under the zero-term convention regardless of the comparison series. No one interpretation is silently chosen.

The source supplies a Borwein comparison theorem for the different ratio \(c_n/b_n\): if \(xf(x)\) is positive and nonincreasing for \(x\ge a>0\), and \(\sum_{n\ge1}f(n)<\infty\), then \(\sum_{n\ge1}f(c_n/b_n)<\infty\). This is retained as a supplied theorem statement, not confused with the requested \(c_n/a_n\) case. The supplied OPEN-TRIAGE summary reports open status in the 2018 edition and no definitive resolution located.
\subsection{Short English statement}
Characterize functions preserving convergence when evaluated at the ratio of a twice-summed sequence to its original terms, and compare with the square-index series. Zero terms, the function class, and the intended convergence implication need clarification.
\subsection{Sources}
\begin{itemize}
\item Local UnsolvedMath registry, permanent ID \(2307031\), exact source code \texttt{catalogue entry}, ``Research Problems in Function Theory --- Problem 7.31.''
\item Hayman, \emph{Research Problems in Function Theory} (2018), Problem 7.31, attributed to W. K. Hayman; \url{https://arxiv.org/abs/1809.07200}. Source-TeX extraction; recorded access 2026-07-29.
\item The source attributes the \(c_n/b_n\) comparison theorem to Borwein without a complete reference. The supplied OPEN-TRIAGE summary gives no separate current-resolution citation.
\end{itemize}
\end{document}
